% titel: Blatt 5, A1: Born-Haber-Kreisprozess (Gitterenergie)
% skript: Mortimer Kap. 7 (chemie: Ionenbindung, Gitterenergie)

Die Gitterenergie $U$ einer Ionenverbindung ist die Energie, die bei der
Bildung des Kristalls aus den getrennten Gas-Ionen frei wird. Sie ist nicht
direkt messbar, lässt sich aber über einen Kreisprozess (Born-Haber) aus
lauter einzeln messbaren Enthalpien bestimmen.

\begin{werkzeugkasten}[Born-Haber-Kreisprozess]
Für ein Salz $MX_n$ aus Metall $M$ und Nichtmetall $X_2$ gilt nach Hess:
\begin{equation}
  \Delta H_{\mathrm f} =
  \Delta H_{\mathrm{sub}} + \sum \mathrm{IE} + \Delta H_{\mathrm{d}} + \sum \mathrm{EA} + U .
\end{equation}
Dabei bedeuten: $\Delta H_{\mathrm{sub}}$ Sublimationsenthalpie des Metalls
$M(\mathrm s)\to M(\mathrm g)$; $\sum\mathrm{IE}$ die Summe aller benötigten
Ionisierungsenergien zum Aufbau $M\to M^{n+}$; $\Delta H_{\mathrm d}$ Verdampfung
und Dissoziation des Nichtmetalls (z.\ B. $\mathrm{Br_2(l)}\to 2\,\mathrm{Br(g)}$,
beide Schritte in einem Wert zusammengefasst); $\sum\mathrm{EA}$ die Summe der
Elektronenaffinitäten der $n$ Nichtmetall-Atome; $U$ die gesuchte Gitterenergie.
Nach $U$ aufgelöst:
\begin{equation}
  U = \Delta H_{\mathrm f} -
      \big(\Delta H_{\mathrm{sub}} + \sum\mathrm{IE} + \Delta H_{\mathrm d} + \sum\mathrm{EA}\big) .
\end{equation}
\textit{Skript:} Mortimer Kap. 7 (Ionenbindung, Born-Haber).
\end{werkzeugkasten}

\begin{achtung}[Darauf achten]
\begin{itemize}
  \item Vorzeichen: Elektronenaffinitäten sind negativ (Energie wird frei),
        Ionisierungsenergien und Sublimation positiv; konsequent durchziehen.
  \item Pro Nichtmetall-Atom genau eine Elektronenaffinität --- bei zwei Atomen
        im Dissoziationsprodukt also der $-2$-fache Wert.
  \item $\Delta H_{\mathrm d}$ liefert bereits die zwei Gas-Atome
        ($X_2\to 2X$); $U$ wird negativ (Kristallbildung ist exotherm).
\end{itemize}
\end{achtung}
